\chapter{Kriterienkatalog zur Plattformbewertung}

\section{Kriterienkatalog für die Plattformbewertung}
\label{sec:kriterienkatalog}

Um im Vorprojekt eine nachvollziehbare und später begründbare Vorauswahl treffen zu können,
wird ein Kriterienkatalog definiert, anhand dessen die betrachteten Open-Source-Cyber-Ranges
strukturiert beschrieben und bewertet werden. Der Katalog ist so gewählt, dass er die späteren
Ziele der Masterarbeit abbildet, ohne dabei bereits eine bestimmte Plattform vorauszusetzen.
Die Kriterien dienen im weiteren Verlauf als Grundlage für die vergleichende Analyse und für
eine begründete Empfehlung.

\textbf{K1 -- Reproduzierbarkeit und Wiederholbarkeit:}
Kann die Plattform identische Umgebungen und Szenarien in gleicher Konfiguration wiederholt bereitstellen
(z.\,B.\ durch versionierbare Definitionen und deterministische Provisionierung)?

\textbf{K2 -- Automatisierung und Lifecycle-Steuerung:}
Unterstützt die Plattform den Lifecycle von Umgebungen (Provisionierung, Start/Stop, Reset, Cleanup)
weitgehend automatisiert und mit geringem manuellem Betriebsaufwand?

\textbf{K3 -- Integrationsfähigkeit externer Werkzeuge:}
Erlaubt die Plattform die Einbindung zusätzlicher Tools innerhalb der Sandbox oder am Rand der Umgebung
(z.\,B.\ Sensoren, Agents, SIEM/Logging, Monitoring), sodass Vergleiche unter gleichen Bedingungen möglich sind?

\textbf{K4 -- Beobachtbarkeit und Datenzugang:}
Stellt die Plattform Mechanismen bereit, um Mess- und Auswertungsdaten systematisch zu erfassen und zugänglich zu machen
(z.\,B.\ Logs, Metriken, Events, Netzwerkdaten) und diese für Analysen zu exportieren bzw.\ weiterzuverarbeiten?

\textbf{K5 -- Zugriffsmöglichkeiten und Nutzerinteraktion:}
Bietet die Plattform praktikable Zugriffsmöglichkeiten auf die bereitgestellten Systeme (z.\,B.\ SSH sowie optional GUI-/Web-Konsole/RDP/VNC)
und unterstützt sie eine nutzerfreundliche Interaktion während der Durchführung?

\textbf{K6 -- Benutzer-, Rollen- und Rechteverwaltung:}
Unterstützt die Plattform Mehrbenutzerbetrieb mit Rollen- und Gruppenkonzepten, nachvollziehbarer Zugriffskontrolle
und klaren Mechanismen für Authentisierung/Autorisierung?

\textbf{K7 -- Wartbarkeit, Dokumentation und Community:}
Ist die Plattform realistisch betreib- und wartbar (z.\,B.\ nachvollziehbarer Installations-/Update-Prozess),
gut dokumentiert und durch aktive Weiterentwicklung bzw.\ Austausch-/Support-Kanäle unterstützt?\newline
\newline


